% Chunk: §3.1 `In' and `Out' States
% Source: Weinberg QFT Vol I, printed pp. 107-112 (PDF pp. 130-135)
% Status: transcribed; source-reviewed; compile-clean; render-reviewed
% Physical pages: 130--135. Figures: none. Uncertainties: none.

The general principles of relativistic quantum mechanics described in the
previous chapter have so far been applied here only to states of a single
stable particle. Such one-particle states by themselves are not very exciting
--- it is only when two or more particles interact with each other that
anything interesting can happen. But experiments do not generally follow the
detailed course of events in particle interactions. Rather, the paradigmatic
experiment (at least in nuclear or elementary particle physics) is one in
which several particles approach each other from a macroscopically large
distance, and interact in a microscopically small region, after which the
products of the interaction travel out again to a macroscopically large
distance. The physical states before and after the collision consist of
particles that are so far apart that they are effectively non-interacting, so
they can be described as direct products of the one-particle states discussed
in the previous chapter. In such an experiment, all that is measured is the
probability distribution, or `cross-sections', for transitions between the
initial and final states of distant and effectively non-interacting particles.
This chapter will outline the formalism~\hyperref[ref:3.1]{[1]} used for calculating these
probabilities and cross-sections.

\subsection{`In' and `Out' States}

A state consisting of several non-interacting particles may be regarded as
one that transforms under the Poincar\'e group as a
direct product of one-particle states. To label the one-particle states we use
their four-momenta $p^\mu$, spin $z$-component (or, for massless particles,
helicity) $\sigma$, and, since we now may be dealing with more than one
species of particle, an additional discrete label $n$ for the particle type,
which includes a specification of its mass, spin, charge, etc. The general
transformation rule is
\begin{align}
U(\Lambda,a)\ket{p_1,\sigma_1,n_1;\,p_2,\sigma_2,n_2;\,\ldots}
={}&\exp\!\left[-ia_\mu\bigl(p_1^\mu+p_2^\mu+\cdots\bigr)\right]
\sqrt{\frac{(\Lambda p_1)^0(\Lambda p_2)^0\cdots}{p_1^0p_2^0\cdots}}
\notag\\
&\cdot\sum_{\sigma'_1\sigma'_2\cdots}
D^{(j_1)}_{\sigma'_1\sigma_1}\bigl(W(\Lambda,p_1)\bigr)
D^{(j_2)}_{\sigma'_2\sigma_2}\bigl(W(\Lambda,p_2)\bigr)\cdots
\notag\\
&\cdot\ket{\Lambda p_1,\sigma'_1,n_1;\,
             \Lambda p_2,\sigma'_2,n_2;\,\ldots}.
\tag{3.1.1}\label{eq:3.1.1}
\end{align}
where $W(\Lambda,p)$ is the Wigner rotation \eqref{eq:2.5.10}, and
$D^{(j)}_{\sigma'\sigma}(W)$ are the conventional $(2j+1)$-dimensional
unitary matrices representing the three-dimensional rotation group. (This is
for massive particles; for any massless particle, the matrix
$D^{(j)}_{\sigma'\sigma}(W(\Lambda,p))$ is replaced with
$\delta_{\sigma'\sigma}\exp(i\sigma\theta(\Lambda,p))$, where $\theta$ is
the angle defined by Eq.~\eqref{eq:2.5.43}.) The states are normalized as in
Eq.~\eqref{eq:2.5.19}:
\begin{align}
&\braket{p'_1,\sigma'_1,n'_1;\,p'_2,\sigma'_2,n'_2;\,\ldots}
 {p_1,\sigma_1,n_1;\,p_2,\sigma_2,n_2;\,\ldots}
\notag\\
&\quad={}
\delta^3(\mathbf p'_1-\mathbf p_1)\delta_{\sigma'_1\sigma_1}
\delta_{n'_1n_1}
\delta^3(\mathbf p'_2-\mathbf p_2)\delta_{\sigma'_2\sigma_2}
\delta_{n'_2n_2}\cdots
\notag\\
&\qquad\mathrel{\phantom{={}}}\pm\text{permutations}.
\tag{3.1.2}\label{eq:3.1.2}
\end{align}
with the term `$\pm$ permutations' included to take account of the
possibility that it is some permutation of the particle types
$n'_1,n'_2,\ldots$ that are of the same species of the particle types
$n_1,n_2,\ldots$. (As discussed more fully in Chapter~4, its sign is $-1$
if this permutation includes an odd permutation of half-integer spin
particles, and otherwise $+1$. This will not be important in the work of the
present chapter.)

We often use an abbreviated notation, letting one Greek letter, say $\alpha$,
stand for the whole collection
$p_1,\sigma_1,n_1;\,p_2,\sigma_2,n_2;\,\ldots$. In this notation,
Eq.~\eqref{eq:3.1.2} is written simply
\begin{equation}
\braket{\alpha'}{\alpha}=\delta(\alpha'-\alpha),
\tag{3.1.3}\label{eq:3.1.3}
\end{equation}
with $\delta(\alpha'-\alpha)$ standing for the sum of products of delta
functions and Kronecker deltas appearing on the right-hand side of
Eq.~\eqref{eq:3.1.2}. Also, in summing over states, we write
\begin{equation}
\int d\alpha\,\cdots\equiv
\sum_{n_1\sigma_1n_2\sigma_2\cdots}
\int d^3p_1\,d^3p_2\cdots.
\tag{3.1.4}\label{eq:3.1.4}
\end{equation}
In particular, the completeness relation for states normalized as in
Eq.~\eqref{eq:3.1.3} reads
\begin{equation}
\ket{\Psi}=\int d\alpha\,\ket{\alpha}\braket{\alpha}{\Psi}.
\tag{3.1.5}\label{eq:3.1.5}
\end{equation}

The transformation rule \eqref{eq:3.1.1} is only possible for particles that for one
reason or another are not interacting. Setting
$\Lambda^\mu{}_{\nu}=\delta^\mu{}_{\nu}$ and $a^\mu=(\tau,0,0,0)$, for
which $U(\Lambda,a)=\exp(iH\tau)$, Eq.~\eqref{eq:3.1.1} requires among other things
that $\ket{\alpha}$ be an energy eigenstate,
\begin{equation}
H\ket{\alpha}=E_\alpha\ket{\alpha},
\tag{3.1.6}\label{eq:3.1.6}
\end{equation}
with an energy equal to the sum of the one-particle energies,
\begin{equation}
E_\alpha=p_1^0+p_2^0+\cdots,
\tag{3.1.7}\label{eq:3.1.7}
\end{equation}
and with no interaction terms, terms that would involve more than one
particle at a time.

On the other hand, the transformation rule \eqref{eq:3.1.1} does apply in scattering
processes at times $t\to\mp\infty$. As explained at the beginning of this
chapter, in the typical scattering experiment we start with particles at time
$t\to-\infty$ so far apart that they are not yet interacting, and end with
particles at $t\to+\infty$ so far apart that they have ceased interacting. We
therefore have not one but two sets of states that transform as in
Eq.~\eqref{eq:3.1.1}: the `in' and `out' states\footnote{The source uses the
labels `$+$' and `$-$' for `in' and `out' states, respectively. This may seem
backward, but the convention has become traditional; it arises from the signs
in Eq.~\eqref{eq:3.1.16}. This edition writes the labels explicitly as
`$\mathrm{in}$' and `$\mathrm{out}$'.}
$\InKet{\alpha}$ and $\OutKet{\alpha}$
will be found to contain the particles described by the label $\alpha$ if
observations are made at $t\to-\infty$ or $t\to+\infty$,
respectively.

Note how this definition is framed. To maintain manifest Lorentz invariance,
in the formalism we are using here, state-vectors do not change with time ---
a state-vector $\ket{\Psi}$ describes the whole spacetime history of a system
of particles. (This is known as the \textit{Heisenberg picture}, in
distinction with the Schr\"odinger picture, where the operators are constant
and the states change with time.) Thus we do \textit{not} say that
$\InKet{\alpha}$ and $\OutKet{\alpha}$ are the
limits at $t\to-\infty$ and $t\to+\infty$ of a time-dependent state-vector
$\ket{\Psi(t)}$.

However, implicit in the definition of the states is a choice of the inertial
frame from which the observer views the system; different observers see
\textit{equivalent} state-vectors, but not the \textit{same} state-vector. In
particular, suppose that a standard observer $O$ sets his or her clock so that
$t=0$ is at some time during the collision process, while some other observer
$O'$ at rest with respect to the first uses a clock set so that $t'=0$ is at a
time $t=\tau$; that is, the two observers' time coordinates are related by
$t'=t-\tau$. Then if $O$ sees the system to be in a state $\ket{\Psi}$, $O'$
will see the system in a state
$U(\mathbf{1},-\tau)\ket{\Psi}=\exp(-iH\tau)\ket{\Psi}$. Thus the appearance
of the state long before or long after the collision (in whatever basis is
used by $O$) is found by applying a time-translation operator
$\exp(-iH\tau)$ with $\tau\to-\infty$ or $\tau\to+\infty$, respectively. Of
course, if the state is really an energy eigenstate, then it cannot be
localized in time --- the operator $\exp(-iH\tau)$ yields an inconsequential
phase factor $\exp(-iE_\alpha\tau)$. Therefore, we must consider
wave-packets, superpositions $\int d\alpha\,g(\alpha)\ket{\alpha}$ of states,
with an amplitude $g(\alpha)$ that is non-zero and smoothly varying over some
finite range $\Delta E$ of energies. The `in' and `out' states are defined so
that the superposition
\[
\exp(-iH\tau)\int d\alpha\,g(\alpha)
\InOutKet{\alpha}
=\int d\alpha\,e^{-iE_\alpha\tau}g(\alpha)
\InOutKet{\alpha}
\]
has the appearance of a corresponding superposition of free-particle states
for $\tau\ll-1/\Delta E$ or $\tau\gg+1/\Delta E$, respectively.

To make this concrete, suppose we can divide the time-translation generator
$H$ into two terms, a free-particle Hamiltonian $H_0$ and an interaction
$H_{\mathrm{int}}$,
\begin{equation}
H=H_0+H_{\mathrm{int}},
\tag{3.1.8}\label{eq:3.1.8}
\end{equation}
in such a way that $H_0$ has eigenstates $\ket{\alpha}$ that have the same
appearance as the eigenstates $\InKet{\alpha}$ and
$\OutKet{\alpha}$ of the complete Hamiltonian,
\begin{equation}
H_0\ket{\alpha}=E_\alpha\ket{\alpha},
\tag{3.1.9}\label{eq:3.1.9}
\end{equation}
\begin{equation}
\braket{\alpha'}{\alpha}=\delta(\alpha'-\alpha).
\tag{3.1.10}\label{eq:3.1.10}
\end{equation}
Note that $H_0$ is assumed here to have the same spectrum as the full
Hamiltonian $H$. This requires that the masses appearing in $H_0$ be the
physical masses that are actually measured, which are not necessarily the
same as the `bare' mass terms appearing in $H$; the difference if there is any
must be included in the interaction $H_{\mathrm{int}}$, not $H_0$. Also, any
relevant bound states in the spectrum of $H$ should be introduced into $H_0$
as if they were elementary particles.\footnote{Alternatively, in
non-relativistic problems we can include the binding potential in $H_0$. In
the application of this method to `rearrangement collisions,' where some bound
states appear in the initial state but not the final state, or vice-versa, one
must use a different split of $H$ into $H_0$ and $H_{\mathrm{int}}$ in the
initial and final states.}

The `in' and `out' states can now be defined as eigenstates of $H$, not $H_0$,
\begin{equation}
H\InOutKet{\alpha}
=E_\alpha\InOutKet{\alpha},
\tag{3.1.11}\label{eq:3.1.11}
\end{equation}
which satisfy the condition
\begin{align}
\int d\alpha\,e^{-iE_\alpha\tau}g(\alpha)
\InKet{\alpha}
&\longrightarrow
\int d\alpha\,e^{-iE_\alpha\tau}g(\alpha)\ket{\alpha}
&& (\tau\to-\infty),
\notag\\
\int d\alpha\,e^{-iE_\alpha\tau}g(\alpha)
\OutKet{\alpha}
&\longrightarrow
\int d\alpha\,e^{-iE_\alpha\tau}g(\alpha)\ket{\alpha}
&& (\tau\to+\infty).
\tag{3.1.12}\label{eq:3.1.12}
\end{align}
Eq.~\eqref{eq:3.1.12} can be rewritten as the requirement that:
\begin{align*}
\exp(-iH\tau)\int d\alpha\,g(\alpha)
\InOutKet{\alpha}
\longrightarrow
\exp(-iH_0\tau)\int d\alpha\,g(\alpha)\ket{\alpha}
\end{align*}
for $\tau\to-\infty$ or $\tau\to+\infty$, respectively. This is sometimes
rewritten as a formula for the `in' and `out' states:
\begin{align}
\InKet{\alpha}&=\Omega(-\infty)\ket{\alpha},
\notag\\
\OutKet{\alpha}&=\Omega(+\infty)\ket{\alpha},
\tag{3.1.13}\label{eq:3.1.13}
\end{align}
where
\begin{equation}
\Omega(\tau)\equiv\exp(+iH\tau)\exp(-iH_0\tau).
\tag{3.1.14}\label{eq:3.1.14}
\end{equation}
However, it should be kept in mind that $\Omega(\mp\infty)$ in
Eq.~\eqref{eq:3.1.13} gives meaningful results only when acting on a smooth
superposition of energy eigenstates.

One immediate consequence of the definition \eqref{eq:3.1.12} is that the `in' and
`out' states are normalized just like the free-particle states. To see this,
note that since the left-hand side of Eq.~\eqref{eq:3.1.12} is obtained by letting the
unitary operator $\exp(-iH\tau)$ act on a time-independent state, its norm is
independent of time, and therefore equals the norm of its limit for
$\tau\to\mp\infty$, i.e., the norm of the right-hand side of Eq.~\eqref{eq:3.1.12}:
\begin{align*}
&\int d\alpha\,d\beta\,
\exp\bigl(-i(E_\alpha-E_\beta)\tau\bigr)
g(\alpha)g^*(\beta)
\InOutBra{\beta}\InOutKet{\alpha}
\\
&\qquad={}
\int d\alpha\,d\beta\,
\exp\bigl(-i(E_\alpha-E_\beta)\tau\bigr)
g(\alpha)g^*(\beta)\braket{\beta}{\alpha}.
\end{align*}
Since this is supposed to be true for all smooth functions $g(\alpha)$, the
scalar products must be equal:
\begin{equation}
\InOutBra{\beta}\InOutKet{\alpha}
=\braket{\beta}{\alpha}=\delta(\beta-\alpha).
\tag{3.1.15}\label{eq:3.1.15}
\end{equation}

It is useful for some purposes to have an explicit though formal solution of
the energy eigenvalue equation \eqref{eq:3.1.11} satisfying the conditions \eqref{eq:3.1.12}.
For this purpose, write Eq.~\eqref{eq:3.1.11} as
\[
(E_\alpha-H_0)\InOutKet{\alpha}
=H_{\mathrm{int}}\InOutKet{\alpha}.
\]
The operator $E_\alpha-H_0$ is not invertible; it annihilates not only the
free-particle state $\ket{\alpha}$, but also the continuum of other
free-particle states $\ket{\beta}$ of the same energy. Since the `in' and
`out' states become just $\ket{\alpha}$ for $H_{\mathrm{int}}\to0$, we
tentatively write the formal solutions as $\ket{\alpha}$ plus a term
proportional to $H_{\mathrm{int}}$:
\begin{align}
\InKet{\alpha}
&=\ket{\alpha}+(E_\alpha-H_0+i\epsilon)^{-1}
H_{\mathrm{int}}\InKet{\alpha},
\notag\\
\OutKet{\alpha}
&=\ket{\alpha}+(E_\alpha-H_0-i\epsilon)^{-1}
H_{\mathrm{int}}\OutKet{\alpha},
\tag{3.1.16}\label{eq:3.1.16}
\end{align}
or, expanding in a complete set of free-particle states,
\begin{align}
\InKet{\alpha}
&=\ket{\alpha}+\int d\beta\,
\frac{T^{\mathrm{in}}_{\beta\alpha}\ket{\beta}}
{E_\alpha-E_\beta+i\epsilon},
\notag\\
\OutKet{\alpha}
&=\ket{\alpha}+\int d\beta\,
\frac{T^{\mathrm{out}}_{\beta\alpha}\ket{\beta}}
{E_\alpha-E_\beta-i\epsilon},
\tag{3.1.17}\label{eq:3.1.17}
\end{align}
\begin{equation}
T^{\mathrm{in/out}}_{\beta\alpha}
\equiv\bra{\beta}H_{\mathrm{int}}\InOutKet{\alpha},
\tag{3.1.18}\label{eq:3.1.18}
\end{equation}
with $\epsilon$ a positive infinitesimal quantity, inserted to give meaning to
the reciprocal of $E_\alpha-H_0$. These are known as the
\textit{Lippmann--Schwinger equations}~\hyperref[ref:3.1a]{[1a]}. We shall use Eq.~\eqref{eq:3.1.17} at the
end of the next section to give a slightly less unrigorous proof of the
orthonormality of the `in' and `out' states.

It remains to be shown that Eq.~\eqref{eq:3.1.17}, with a $+i\epsilon$ or
$-i\epsilon$ in the denominator, satisfies the condition \eqref{eq:3.1.12} for an
`in' or an `out' state, respectively. For this purpose, consider the
superpositions
\begin{equation}
\InOutKet{\Psi_g(t)}
\equiv\int d\alpha\,e^{-iE_\alpha t}g(\alpha)
\InOutKet{\alpha},
\tag{3.1.19}\label{eq:3.1.19}
\end{equation}
\begin{equation}
\ket{\Phi_g(t)}\equiv
\int d\alpha\,e^{-iE_\alpha t}g(\alpha)\ket{\alpha}.
\tag{3.1.20}\label{eq:3.1.20}
\end{equation}
We want to show that $\InKet{\Psi_g(t)}$ and
$\OutKet{\Psi_g(t)}$ approach $\ket{\Phi_g(t)}$ for
$t\to-\infty$ and $t\to+\infty$, respectively. Using Eq.~\eqref{eq:3.1.17} in
Eq.~\eqref{eq:3.1.19} gives
\begin{align}
\InOutKet{\Psi_g(t)}
=\ket{\Phi_g(t)}+\int d\alpha\int d\beta\,
\frac{e^{-iE_\alpha t}g(\alpha)
T^{\mathrm{in/out}}_{\beta\alpha}\ket{\beta}}
{E_\alpha-E_\beta\mathbin{\pm}i\epsilon},
\tag{3.1.21}\label{eq:3.1.21}
\end{align}
where the upper and lower signs refer to `in' and `out', respectively. Let us
recklessly interchange the order of integration, and consider first the
integrals
\[
\mathcal{I}^{\mathrm{in/out}}_\beta\equiv
\int d\alpha\,
\frac{e^{-iE_\alpha t}g(\alpha)
T^{\mathrm{in/out}}_{\beta\alpha}}
{E_\alpha-E_\beta\mathbin{\pm}i\epsilon}.
\]
For $t\to-\infty$, we can close the contour of integration for the energy
variable $E_\alpha$ in the upper half-plane with a large semi-circle, with the
contribution from this semi-circle killed by the factor
$\exp(-iE_\alpha t)$, which is exponentially small for $t\to-\infty$ and
$\operatorname{Im}E_\alpha>0$. The integral is then given by a sum over the
singularities of the integrand in the upper half-plane. The functions
$g(\alpha)$ and $T^{\mathrm{in/out}}_{\beta\alpha}$ may, in general, be
expected to have some singularities at values of $E_\alpha$ with finite
positive imaginary parts, but just as for the large semi-circle, their
contribution is exponentially damped for $t\to-\infty$. (Specifically, $-t$
must be much greater than both the time-uncertainty in the wave-packet
$g(\alpha)$ and the duration of the collision, which respectively govern the
location of the singularities of $g(\alpha)$ and
$T^{\mathrm{in/out}}_{\beta\alpha}$ in the complex $E_\alpha$ plane.) This
leaves the singularity in
$(E_\alpha-E_\beta\mathbin{\pm}i\epsilon)^{-1}$, which is in the upper
half-plane for $\mathcal{I}^{\mathrm{out}}_\beta$ but not
$\mathcal{I}^{\mathrm{in}}_\beta$. We conclude then that
$\mathcal{I}^{\mathrm{in}}_\beta$ vanishes for $t\to-\infty$. In the same
way, for $t\to+\infty$ we must close the contour of integration in the lower
half-plane, and so $\mathcal{I}^{\mathrm{out}}_\beta$ vanishes in this limit.
We conclude that $\InOutKet{\Psi_g(t)}$ approaches
$\ket{\Phi_g(t)}$ for $t\to\mp\infty$, in agreement with the defining
condition \eqref{eq:3.1.12}.

\begin{center}
$*\ *\ *$
\end{center}

For future use, we note a convenient representation of the factor
$(E_\alpha-E_\beta\mathbin{\pm}i\epsilon)^{-1}$ in Eq.~\eqref{eq:3.1.17}. In
general, we can write
\begin{equation}
(E\mathbin{\pm}i\epsilon)^{-1}
=\frac{\Pi_\epsilon}{E}\mp i\pi\delta_\epsilon(E),
\tag{3.1.22}\label{eq:3.1.22}
\end{equation}
